\documentclass[a4paper]{article}
\usepackage{standalone}
\usepackage{import}
\input{../../../../newpreamble.tex}
\title{Math 299 Exercises}
\author{Lance Remigio}

\begin{document}
\maketitle    
\tableofcontents

\section{Week 1}

\import{week}{week1/week1.tex}

Since \( h: \R^{n} \to \R^{n}  \) is an isometry, we know from definition that 
\[  d(h(x), h(y)) = d(x,y) \]
for all \( x,y \in \R^{n} \). In particular, since the norm on \( \R^{n} \) induces the distance function defined by
\[  d(x,y) = \|x - y\| \]
we also have that the norm is preserved. Indeed, we have for all \( x,y \in \R^{n} \),
\[  \|h(x) - h(y)\| = d(h(x), h(y)) = d(x,y) = \|x-y\|. \]
Since \( h(0) = 0  \) (by assumption), it is the case that the norm is also preserved. Indeed, for all \( x\in \R^{n} \)
, we have 
\[  \|h(x)\| = \|h(x) - h(0)\| = \|x - 0\|  = \|x\|.\]
The fact that the norm is preserved, it also tells us that the inner product \( \langle \cdot , \cdot \rangle \) on \( \R^{n} \) is preserved. That is, if for all \( x \in \R^{n} \) and \( \langle \cdot , \cdot \rangle = \|\cdot\|^{2}  \), then we have 
\begin{align*}
    \langle h(x), h(x) \rangle &= \|h(x) \|^{2} \\
                                              &= \|x\|^{2} \\
                                              & \langle x  ,x \rangle.
\end{align*}
 

\end{document}

